\title{LongRoPE: Extending LLM Context Window Beyond 2 Million Tokens}

\begin{abstract}
A substantial context window is an advantageous attribute for large language models (LLMs). Nevertheless, the existing scope of extended context windows remains constrained to approximately 128k tokens. This limitation arises from factors such as significant expenses associated with fine-tuning, a limited supply of lengthy textual data, and disruptive value shifts that occur when incorporating novel token positions.

This study presents LongRoPE, which, for the first time, expands the context window of pre-trained LLMs to a remarkable \textbf{2048k} tokens, requiring at most 1k fine-tuning steps on sequences up to 256k tokens in length, while preserving the model's capabilities on its original short context. This advancement relies on three central contributions: (i) we uncover and leverage two distinct types of non-uniformities present in positional interpolation through an efficient search algorithm, yielding improved fine-tuning initialization and facilitating an 8$\times$ extension in zero-shot settings; (ii) we propose a gradual extension technique, beginning with the fine-tuning of a 256k-token LLM, followed by an additional round of positional interpolation on the already extended model to reach a 2048k context window; (iii) we refine LongRoPE on sequences of length 8k to restore performance for short context windows. Comprehensive evaluations on LLaMA2 and Mistral across a variety of benchmarks highlight the efficacy of our method. LLMs augmented with LongRoPE retain their original architectural structure, requiring only minimal changes to the positional embedding mechanism, and remain compatible with the vast majority of previous optimizations. Code will be made available at \texttt{https://github.com/microsoft/LongRoPE}

\section{Introduction}

Although Large Language Models (LLMs) have achieved impressive results across a range of tasks~\cite{gpt4,llama2}, their capabilities are constrained by a relatively small context window, such as the 4096-token boundary found in LLaMA2~\cite{llama2}. When input exceeds this limit, model performance deteriorates, largely because these extended positions fall outside of the data distribution the model encountered during training. This limitation creates difficulties for applications that demand the processing of extensive contextual information, including scenarios like in-context learning with multiple examples~\cite{Huang2023FewerIM} and tasks involving LLM agents~\cite{park2023generative,madaan2023selfrefine}.

Recent studies have demonstrated that the context window of a pre-trained LLM can be expanded to approximately 128k tokens by fine-tuning with lengthier documents~\cite{longlora,pi,yarn,zhang2024soaring,liu2023scaling}. However, scaling the context window further faces three significant challenges. Firstly, the addition of previously unused position indices during extension results in a high frequency of catastrophic values, causing distributional shifts that make the fine-tuning process unstable and slow to converge~\cite{pi}. This issue is especially pronounced in cases where an increase from 4k to $>$1000k token positions means that over 90\% of positions are uninitialized. Secondly, effective fine-tuning generally necessitates access to documents matching the target extended lengths, yet contemporary datasets seldom include texts beyond 1000k tokens. Additionally, processing these exceptionally long documents is both computationally demanding and requires a substantial investment in training time and GPU capacity. Thirdly, as the context window grows to extreme lengths, the model’s attention becomes increasingly diluted, distributed over a vast number of token positions. This leads to diminished performance on tasks requiring comprehension over the model’s original, shorter context lengths~\cite{pi}.

A common strategy to address the initial issue is to apply interpolation to RoPE positional embeddings~\cite{rope,pi}, which rescales position indices beyond the pretrained limit back into the original range, as depicted in Fig.\ref{fig:overview}. Position Interpolation (PI)~\cite{pi} conducts a linear scaling of RoPE’s rotary angles according to the extension ratio. The NTK method~\cite{ntk,dynamicntk} introduces non-uniform interpolation and extrapolation along different RoPE dimensions. Meanwhile, YaRN~\cite{yarn} separates RoPE dimensions into three groups based on their frequency and employs a mixture of extrapolation, NTK-based, and linear interpolation methods for each group. Nonetheless, within the Transformer framework, positional embeddings possess a \emph{complex, non-uniform information entropy}. Current techniques insufficiently utilize this nuanced non-uniformity, resulting in loss of information and consequently restricting the effective extension of the context window.

Section~\ref{sec:analysis} empirically demonstrates two principal observations: \textbf{(1)} Successful positional interpolation requires accounting for two types of non-uniformity—differences across RoPE dimensions and among token positions. Reduced interpolation is more advantageous for dimensions at the lower end and for early token positions, though the optimal approach is contingent upon the intended length of the sequence extension. \textbf{(2)} Integrating these non-uniformities within positional interpolation facilitates better preservation of the original RoPE's information, especially in crucial dimensions and early token positions. This strategy reduces the information loss typically associated with positional interpolation, thereby leading to improved initial conditions for subsequent fine-tuning. Additionally, this method enables sequence length extensions up to 8$\times$ in settings without fine-tuning.

Building on these insights, we present \textbf{LongRoPE}, a robust approach designed to expand the LLM context window to beyond 2 \emph{million} tokens. The effectiveness of LongRoPE derives from three central innovations. First, {LongRoPE} harnesses multidimensional non-uniform characteristics inherent in positional interpolation to their fullest potential. This is accomplished by determining optimal rescale factors for the rotation angles in RoPE along each dimension, taking into account token position dependencies. Given that the search space for these rescale factors grows exponentially with higher extension ratios, {LongRoPE} employs an evolutionary search strategy that incorporates two specific optimization methods to significantly improve search efficiency. An illustrative example of the rescaled RoPE generated through this search process is provided in Fig.~\ref{fig:overview}.

Subsequently, {LongRoPE} implements a resource-efficient, stepwise expansion method to attain a 2048k context window, circumventing the necessity for explicit fine-tuning on extremely lengthy sequences, which are rarely available. Initially, the approach identifies a suitable 256k context length for the pre-trained LLM and proceeds with fine-tuning at this length. Next, by exploiting our non-uniform positional interpolation—which supports an 8$\times$ length increase even in the absence of further fine-tuning—a second search for appropriate RoPE rescale factors is performed on this fine-tuned, length-extended LLM. This procedure ultimately enables support for a 2048k context window on both LLaMA2 and Mistral~\cite{mistral}.

In order to prevent a drop in performance when processing shorter context windows, {LongRoPE} also modifies the RoPE rescale factors on the expanded LLM. Following the approach used for scaling up from 256k to 2048k, we apply our search algorithm to decrease the context window size to 4k and 8k on the LLM that has been fine-tuned at 256k, thereby reducing positional interpolation. When the input sequence is shorter than 8k during inference, RoPE is refreshed using the rescale factors determined through our search process.

A comprehensive series of experiments involving multiple LLMs and a range of long-context evaluation tasks validates the efficacy of our approach. Our results indicate that LongRoPE consistently sustains low perplexity across evaluation lengths spanning from 4k up to 2048k, attains passkey retrieval accuracy exceeding 90\%, and achieves accuracy on par with standard benchmarks situated within the 4096-token context window. Furthermore, LongRoPE is compatible with any LLM utilizing RoPE embedding. We intend to make both our codebase and the LongRoPE-2048k models publicly available.

\section{Non-uniformity in Positional Interpolation}
\label{sec:analysis}

\subsection{Preliminary}

Transformer architectures depend on the inclusion of explicit positional data—commonly incorporated via position embeddings—to capture the sequential structure of input tokens. In this study, we investigate RoPE~\cite{rope} position embeddings, a method extensively adopted in modern LLMs. Given a token located at position index $n$, its RoPE encoding is succinctly represented by:

\begin{equation}
		\fontsize{7.8}{7.8} \selectfont
	\label{eq:rope}
	\left[ cos(n\theta_0), sin(n\theta_0), cos(n\theta_1), \cdot\cdot\cdot, cos(n\theta_{d/2-1}), sin(n\theta_{d/2-1}) \right]   
\end{equation}

Here, $d$ refers to the embedding dimension, $n\theta_i$ denotes the rotary angle for the token located at position $n$, and $\theta_i=\theta^{-2i/d}$ specifies the rotation frequency. The standard base value chosen for $\theta$ in RoPE is 10000.

\textbf{\emph{Context window extension ratio $s$} and positional interpolation}. The parameter $s$ is defined as the proportion between the extended context window length $L'$ and the initial length $L$: $s=\frac{L'}{L}$.

In order to increase the context window from $L$ to $L'$, existing positional interpolation approaches propose adjusting the rotation frequencies $\theta_i$ downward in accordance with the extension ratio $s$. Defining $\beta = \theta^{2/d}$ and letting $\lambda$ represent the rescaling factor associated with $s$, we can formulate a unified expression for these positional interpolation techniques as follows:

\begin{equation}	
	\fontsize{7.5}{7.5} \selectfont
	\label{eq:unified_rope}
	\begin{aligned}
		\left[cos\left(\frac{n}{\lambda(\beta)^0}\right), sin \left(\frac{n}{\lambda(\beta)^0}\right), cos\left(\frac{n}{\lambda(\beta)^1}\right), \cdot\cdot\cdot,  sin \left(\frac{n}{\lambda(\beta)^{d/2-1}}\right) \right]   
	\end{aligned}
\end{equation}

\textbf{Linear positional interpolation (PI)}. PI~\cite{pi} proposes interpolating position indices linearly when extending sequence length, as long as the new length remains within the range supported by pre-training. Given a desired extension ratio $s$, all positional rotation angles are uniformly scaled down by a factor of $\lambda = s$ for every RoPE dimension. This approach compresses the representation of positional information, which reduces the model's capacity to discriminate between tokens located near each other. As a consequence, PI generally yields suboptimal results when applied to large extension ratios.

\textbf{NTK-based interpolation and extrapolation}. ~\cite{ntk,dynamicntk} analyze RoPE through the lens of information encoding and employ Neural Tangent Kernel (NTK) theory~\cite{jacot2018neural,tancik2020fourier} in their approach. To address the problem of position crowding present in PI, their method redistributes interpolation burden across the various RoPE dimensions. Specifically, dimensions with lower indices (which represent higher frequencies) are scaled less, while those with higher indices (corresponding to lower frequencies) are scaled more significantly, enabling both interpolation and extrapolation over positions, where $\lambda=s^{i}$. The enhanced dynamic NTK~\cite{dynamicntk} further refines this by modulating the extension ratio for each position according to the sequence's actual length. Unlike PI, which requires model fine-tuning, NTK-aware strategies are capable of enlarging the context window in scenarios where fine-tuning is not possible, though they are generally restricted to a maximum extension ratio of 4$\times$.

\textbf{YaRN}~\cite{yarn} presents a notable advancement in enhancing the efficacy of positional interpolation. This method categorizes RoPE dimensions into three separate groups based on their associated frequencies, assigning distinct interpolation approaches to each. Dimensions with high frequencies are handled through extrapolation ($\lambda$=1), whereas dimensions with low frequencies utilize linear interpolation (PI). Those RoPE dimensions with intermediate frequencies are treated using NTK. A central feature of YaRN is its approach to grouping the RoPE dimensions; however, this categorization currently relies on manual, empirically-driven decisions. As a consequence, such reliance on human heuristics may lead to less-than-optimal outcomes when applied to newly developed LLMs.

\subsection{Study on Non-uniform Positional Interpolation}

Drawing from the ideas underlying NTK and YaRN, we observe that their improvements stem from the introduction of non-linear transformations, particularly through addressing distinct frequency patterns across RoPE dimensions for targeted interpolation and extrapolation. Yet, existing approaches to non-linearity predominantly depend on handcrafted heuristics. This leads us to two important inquiries: (1) Does the present method of positional interpolation represent the best possible solution? (2) Might there exist alternative, previously unexamined, forms of non-linearity?

In order to address these questions, we employ evolution search (refer to Sec.~\ref{sec:method}) to identify improved non-uniform positional interpolations for LLaMA2-7B. This search process is steered by perplexity measurements, utilizing five randomly chosen instances from the PG19~\cite{pg19} validation set. Our experimental results lead to several important observations.

\textbf{Finding 1}: \textit{RoPE dimensions exhibit substantial non-uniformities, which are not effectively handled by current positional interpolation methods.}

We determine the optimal value of $\lambda$ for each RoPE dimension as defined in Eq.~\ref{eq:unified_rope}. In Table~\ref{tbl:exp1}, we report the perplexity results of LLaMA2-7B across various approaches on the PG19 and Proof-pile~\cite{proof-pile} test datasets, evaluated without fine-tuning. The configuration identified by our search yields marked perplexity reductions, indicating that prevalent strategies such as linear (PI) and non-uniform (Dynamic-NTK and YaRN) interpolation do not achieve optimal performance. Interestingly, on PG19, YaRN lags behind both PI and NTK, likely because it fails to attain the desired context window length when the LLM is not fine-tuned. Specifically, for an 8k context window, YaRN demonstrates a notable perplexity increase beyond 7k tokens.

Our investigation reveals that the rescaled factors $\lambda$ in Eq.~\ref{eq:unified_rope} exhibit a non-uniform pattern, unlike the constant scale $s$ employed in the formulas of PI, NTK, or the group-based calculation in YaRN. Adopting these non-uniform scaling factors yields substantial gains in LLaMA2's language modeling abilities (specifically, perplexity) across context lengths of 8k and 16k tokens, all without requiring additional fine-tuning. The effectiveness of this approach stems from its capacity to closely maintain the characteristics of the original RoPE, particularly in critical dimensions, thereby lessening the model’s challenge of differentiating tokens located in proximate positions.

\textbf{Finding 2}: \textit{RoPE for the initial tokens in the input sequence should be extrapolated with less interpolation.} 

We propose that the RoPE interpolation for the first $\hat{n}$ tokens in input sequences should be minimized. The rationale behind this is that these tokens tend to obtain high attention weights, making them pivotal in the attention mechanism, as noted by Streaming LLM~\cite{streamingllm} and LM-Infinite~\cite{han2023lminfinite}. To evaluate this hypothesis, we expand the context window to lengths of 8k and 16k utilizing PI and NTK methods, while leaving the initial $\hat n$ (0, 2, ..., 256) tokens exempt from interpolation. When $\hat n$ is set to 0, the methods revert to their original form. As demonstrated in Table~\ref{tbl:tokenposition}, two principal findings emerge: \textbf{(1)} excluding position interpolation for the earliest tokens contributes to better model performance; and \textbf{(2)} the optimal value of $\hat n$ varies according to the length of the context window extension.

\textbf{Finding 3}: \textit{Non-uniform positional interpolation effectively extends LLM context window in both fine-tuning and non-fine-tuning settings.} 

Although we have demonstrated that our optimized non-uniform position interpolation method leads to notable improvements in extension performance at 8k and 16k context lengths without additional training, achieving robust results at even longer context windows necessitates fine-tuning. Therefore, we fine-tune LLaMA2-7B equipped with our optimized RoPE configuration to accommodate a 64k context window (see Appendix for further details on experimental settings). As evidenced in Table~\ref{tbl:finetune}, our approach achieves substantial gains over both PI and YaRN, outperforming them both in the pre- and post-fine-tuning stages on LLaMA2-7B. This can be attributed to our method’s capacity to leverage non-uniform positional interpolation, which effectively reduces information loss and serves as a superior starting point for subsequent fine-tuning.

\textbf{Summary}. Our research identifies two critical sources of non-uniformity: the variability across RoPE dimensions and across token positions. Harnessing these sources within positional interpolation strategies enables marked improvements in extending the context capabilities of large language models.

\section{LongRoPE}

\label{sec:method}
Drawing on these observations, we propose LongRoPE. This approach begins by presenting an effective search algorithm designed to leverage both types of non-uniformity. We then apply this method to successfully extend the LLM context window past 2 million tokens.

\subsection{Problem Formulation}

These two forms of non-uniformity greatly expand the potential solution set and add layers of difficulty to the optimization process. To manage this challenge, we recast the optimization of multidimensional, non-uniform position interpolation as a search problem.

Consider a LLM designed for a context window of length $L'$ and tasked with processing extensive input documents $\mathbf{X}$, where every sequence $\mathbf{x} \in \mathbf{X}$ exceeds $L'$ in token count. The initial rotary angle for the $i^{th}$ dimension in the RoPE embedding at token position $n$ is represented as $\frac{n}{\beta^i}$. The optimization objective can therefore be defined as:

\begin{equation}
\begin{gathered}
		\label{eq:problem}

\underset {\mathbf{x}\in\mathbf{X}; \, |\mathbf{x}|\ge L'}{ \operatorname {arg\,min} } \,  \mathcal{L}\, \left( \text{LLM} (\text{RoPE}, \mathbf{X})\right),	\text{where} 
\\
\scalebox{0.93}{
$\underset {\begin{subarray}{l} i=0,\ldots,\frac{d}{2}-1;\\ n\in[ 0,|\mathbf{x}|);\end{subarray}}{\text{RoPE}_(n)}= \bigl[\ldots,\cos\bigl(\mathbb{I}(\hat{\lambda}_{i}, \hat{n})\frac{n}{\beta^i}\bigr),\;\sin\bigl(\mathbb{I}(\hat{\lambda}_{i}, \hat{n})\frac{n}{\beta^i}\bigr),\ldots\bigr]$}\\
	\text{with } \mathbb{I}(\hat{\lambda}_{i},\hat{n})=\begin{cases}
	1&\mbox{ if $n<\hat{n}$}\\
	{\frac{1}{\lambda}_{i}}&\mbox{ if $n\geq\hat{n}$}
\end{cases}
	\end{gathered}
\end{equation}

In this approach, we define a collection of scaling factors, $\mathbb{I}(\hat{\lambda}_{i},\hat{n})$, aimed at addressing both types of non-uniformity. Here, $\hat{\lambda}_i$ captures non-uniformity across RoPE dimensions, while $\hat{n}$ represents variations in token positions. The factor $\mathbb{I}(\hat{\lambda}_{i},\hat{n})$ is utilized to adjust the rotation angle associated with the $i^{th}$ RoPE dimension; specifically, $\hat\lambda_{i}$ serves as the scaling coefficient and $\hat{n}$ as the threshold for token positions. For token positions preceding $\hat{n}$, that is, for the first $\hat{n}$-$1$ tokens, the scaling coefficient $\hat\lambda_{i}$ is inactive, and the unaltered RoPE rotation angle $\frac{n}{\beta^i}$ is applied. When token positions satisfy $n \geq \hat{n}$, the scaling factor comes into effect.

Suppose we are given a desired context window length of $L'$. The task is to determine the optimal set of rescaling coefficients ($\mathbb{I}(\hat{\lambda}_{0},\hat{n})$, $\mathbb{I}(\hat{\lambda}_{1},\hat{n})$, ... , $\mathbb{I}(\hat{\lambda}_{i},\hat{n})$, ...) for each of the RoPE dimensions, from the first through the $d^{th}$. By applying these rescaled factors to the RoPE in the target $\text{LLM}$, the model is able to minimize its next-token prediction loss, $\mathcal{L}$ (which directly relates to perplexity), when evaluated on input sequences $\mathbf{X}$ of token length $L'$.

\subsection{Searching the Non-uniform Position Interpolation}

In order to address the issue outlined in Eq.~\ref{eq:problem}, we present a straightforward yet robust approach that aims to identify the best possible RoPE rescale factors. This method is designed to take full advantage of the complex, multidimensional variations present in positional embedding.

\textbf{Search space}. Our approach constructs an expansive search space that encapsulates the two forms of non-uniformity. The structure of this search space is detailed in Table~\ref{tbl:searchspace}. More precisely, our framework enables the assignment of unique rescale factors to each dimension of the RoPE embedding. For the sake of streamlined search space formulation, we opt to directly search over $\lambda_i$ and $\hat{n}$, rather than working with $\mathbb{I}(\hat{\lambda}_{i},\hat{n})$, noting that $\hat{\lambda}_{i}=1/\lambda_i$. As indicated in Table~\ref{tbl:searchspace}, the search for $\lambda_i$ spans from a lower bound of 1.0 (corresponding to pure extrapolation) up to an upper limit of $s\times1.25$ (representing a broader range of interpolation than that provided by PI), using a granularity of 0.01 in the step size. Here, $s$ denotes the intended context window scaling factor.

The parameter $\hat{n}$ determines how many of the starting token positions preserve the original RoPE embedding, foregoing position interpolation. In practice, we enable $\hat{n}$ to take values from the set \{0, 1, 2, 4, 8, 12, 16, 20, 24, 28, 32, 64, 128, 256\} during the search. Setting $\hat{n}=0$ indicates that the searched rescale factors are applied to every token position.

\textbf{Evolution-based search}. The search space presented in Table~\ref{tbl:searchspace} encompasses a vast array of positional interpolation options, which renders efficient navigation a complex task. As an illustration, an extension factor of $s=4\times$ results in $400^{128/2}\times14$ = $4\times10^{167}$ potential configurations. This combinatorial explosion grows exponentially with increasing extension ratios. To efficiently tackle this challenge, we employ evolution search~\cite{guo2020single} and propose two optimization strategies that substantially accelerate the search process. The overall approach is depicted in Algorithm~\ref{alg:search}.

\textit{Optimized initial population generation}. Rather than selecting all $P$ rescale factors at random for the initial population, our method explicitly includes the RoPE rescale factors for PI, NTK, and YaRN as predefined individuals. The rest of the population, totaling $P-3$ individuals, is produced by applying random mutations to these three rescale factors, each mutation occurring with probability $p$.

\textit{Monotonically non-decreasing constraint}. Once the initial population has been created, we evaluate the perplexity of each candidate using the LLM. This involves applying the associated RoPE rescale parameters to the target LLM and measuring the perplexity on the input $\mathbf{X}$. The best $k$ candidates, ranked by perplexity, are selected as parents for the evolutionary process. Nevertheless, the enormous size of the search space means that simple mutation and crossover steps may frequently yield suboptimal offspring, leading to excessive and avoidable perplexity evaluations. This issue is exacerbated when $L'$ is large, as each perplexity computation requires substantial inference time.

To tackle this issue, we enforce a constraint of non-decreasing monotonicity on the RoPE rescaled factors that are sampled, such that $\lambda_i\le \lambda_{i+1}$. Only those RoPE configurations meeting this requirement are used for evaluating the LLM's perplexity, thereby streamlining the search process considerably. Concretely, this means that $\lambda_i$ must rise monotonically as the RoPE dimension increases, with $i=0,\ldots,63$. The rationale for this monotonicity across dimensions is informed by NTK theory~\cite{jacot2018neural,tancik2020fourier,ntk}, which posits that lower (higher frequency) dimensions benefit from less interpolation (corresponding to a smaller $\lambda_i$), while higher (lower frequency) dimensions are more suited to increased interpolation (thus a larger $\lambda_i$).

\begin{algorithm}[t]
	\small
	\caption{The search algorithm}
	\textbf{Input:} target LLM, input samples $\mathbf{X}$,   population size $P$, mutation size $N_1$, crossover size $N_2$, max iterations $\mathcal{T}$, mutate probability $p$\\

	\begin{algorithmic}[1]
		\label{alg:search}
		\STATE Top-k $\gets$ $\phi$;	
		\STATE $\text{P}_0$ $\gets$ \textit{Initialize\_population\_with\_optimization} ($P$, $\mathbf{X}$, $p$); 
		\FOR{$i = 1$ to $\mathcal{T}$}
		\STATE \textit{Compute\_perplexity} (LLM, $\text{P}_{i-1}$, $\mathbf{X}$);
		\STATE  Top-k $\gets$ \textit{Update\_Topk} (Top-k, $\text{P}_{i-1}$);
		\STATE $\text{P}_{mutation}$ $\gets$ \textit{Mutation\_with\_mono\_constraint} (Top-k, $N_1$, $p$); 
		\STATE $\text{P}_{crossover}$ $\gets$ \textit{Crossover\_with\_mono\_constraint} (Top-k, $N_2$);
		\STATE $\text{P}_i$ $\gets$ $\text{P}_{mutation}$ $\cup$ $\text{P}_{crossover}$ $\cup$ Top-k;
		\ENDFOR
		\STATE Return the member in Top-k achieving the minimum perplexity;
	\end{algorithmic}
\end{algorithm}

\textbf{8$\times$ extension without fine-tuning}. Through our evolutionary search approach, we successfully determine optimal non-uniform RoPE rescale factors that retain crucial dimensions and token positions, thereby reducing information loss due to interpolation. As illustrated in Fig.\ref{fig:non-ft}, this enables our method to expand LLaMA2’s context window from 4k to 32k tokens without the need for any fine-tuning. Conversely, current approaches like PI and non-uniform variants of NTK and YaRN exhibit a sharp increase in perplexity once the context length is doubled.

\subsection{Extending LLM Context Window to 2048K}
\label{sec:extendto2048k}

\textbf{Progressive extension to 2048k}. Here, we present our approach for scaling the context window of pre-trained LLMs beyond the standard 4k tokens, reaching up to 2048k. Our non-uniform positional interpolation technique enables an 8$\times$ expansion without the need for any fine-tuning. However, when pursuing even greater extensions, such as 512$\times$, fine-tuning becomes essential. A potential strategy involves identifying suitable RoPE rescaling factors tailored for the 2048k window, followed by fine-tuning. Nonetheless, this path is hindered by the substantial computational resources required for such extensive training. Furthermore, as indicated by our observations, achieving robust fine-tuning with such a significant increase in context length proves to be highly challenging (refer to Appendix).

Luckily, LongRoPE demonstrates efficacy for both base and fine-tuned extended LLMs. Thus, we propose a streamlined and incremental approach capable of reaching the desired 2048k context window through only 1k fine-tuning iterations, all performed within a 256k sequence length during training.

$\Diamond$ \textit{LongRoPE search for scaling pre-trained LLMs to 256k}. Using LLaMA2 as a case study, we perform a search to achieve target context window lengths of 128k and 256k. This results in extension factors of 32$\times$ and 64$\times$ at this phase.

$\Diamond$ \textit{Fine-tuning to 256k}. Subsequently, we adapt the pre-trained LLM to support a context window of 256k by additional fine-tuning. The process begins with 400 steps of fine-tuning on LLaMA2 using RoPE rescaled factors configured for 128k. Once this stage is completed, the RoPE rescaled factors are updated to accommodate 256k, after which a further 600 steps of fine-tuning are performed on the resulting checkpoint. This sequential approach yields greater efficiency compared to immediately fine-tuning for the full 256k context window.

$\Diamond$ \textit{LongRoPE-enabled scaling of fine-tuned extended LLMs up to 2048k through targeted search}. In the final phase, we conduct an additional search step using the LLM that has already been fine-tuned for 256k sequence length. This process expands the context window dramatically to a size of 2048k, accomplished without any extra fine-tuning procedures. As a result, the model achieves a context length expansion factor of $512\times$.

\textbf{Recovering performance in the original context window}. When the context window is expanded to a very large size of 2048k, a decrease in performance within the initial context window is observed. This phenomenon is recognized in the literature on positional interpolation~\cite{pi}, since this technique compresses the positional embeddings corresponding to the original context window into a much more confined area of the embedding space. This compression has an adverse effect on model performance. Specifically, applying a 512$\times$ extension factor results in substantial crowding of positional embeddings within the initial 4k context window.

To address this issue, we conduct an additional evolutionary search on the extended LLM to fine-tune RoPE rescale factors specifically for shorter context lengths (such as 4k and 8k). The upper bound for the searched $\lambda$ values is lowered, reflecting the reduced need for positional interpolation at these lengths. At inference time, the LLM adaptively modifies the relevant RoPE rescale factors based on the context.

\section{LongRoPE}

\label{sec:method}
Building upon these observations, we propose LongRoPE. This approach initially incorporates an effective search algorithm designed to leverage both forms of non-uniformity. Utilizing this strategy, LongRoPE extends the LLM context window to exceed 2 million tokens.

\subsection{Problem Formulation}

The presence of these two types of non-uniformity significantly enlarges the solution space and adds layers of complexity to the optimization process. To effectively manage this, we reformulate the multidimensional non-uniform position interpolation task as a search problem.

Consider an LLM configured for a context window of size $L'$ and processing extended input sequences $\mathbf{X}$, with each segment $\mathbf{x} \in \mathbf{X}$ having a token length greater than $L'$. The initial rotary angle in the $i^{th}$ dimension of the RoPE embedding for token position $n$ is expressed as $\frac{n}{\beta^i}$. With this setup, the corresponding optimization problem can be articulated as:

\begin{equation}
\begin{gathered}
		\label{eq:problem}

\underset {\mathbf{x}\in\mathbf{X};\, |\mathbf{x}|\ge L'}{ \operatorname {arg\,min} } \,  \mathcal{L} \left( \text{LLM} (\text{RoPE}, \mathbf{X})\right),	\text{with}
\\
\scalebox{0.93}{
$\underset {\begin{subarray}{l} i=0,\ldots,\frac{d}{2}-1;\\ n\in[ 0,|\mathbf{x}|);\end{subarray}}{\text{RoPE}_(n)}= {\left[\ldots,\cos\left(\mathbb{I}(\hat{\lambda}_{i}, \hat{n})\cdot\frac{n}{\beta^i}\right), \sin\left(\mathbb{I}(\hat{\lambda}_{i}, \hat{n})\cdot\frac{n}{\beta^i}\right),\ldots\right] }$}\\
	\text{where} \; \mathbb{I}(\hat{\lambda}_{i},\hat{n})=\begin{cases}
	1&\mbox{\;  $n<\hat{n}$}\\
{\frac{1}{\lambda}_{i}}&\mbox{\; $n\geq\hat{n}$}
\end{cases}
	\end{gathered}
\end{equation}

In this context, we define a collection of scaling factors, $\mathbb{I}(\hat{\lambda}_{i},\hat{n})$, designed to address two types of non-uniformities. The variables $\hat{\lambda_i}$ and $\hat{n}$ represent, respectively, the non-uniform scaling across RoPE dimensions and irregularities in token positions. Concretely, $\mathbb{I}(\hat{\lambda}_{i},\hat{n})$ serves to adjust the rotation angle associated with the $i^{th}$ RoPE dimension, with $\hat\lambda_{i}$ functioning as the scaling parameter and $\hat{n}$ as the position-based threshold. For the initial $\hat{n}-1$ token positions, no adjustment via $\hat\lambda_{i}$ is made, so the standard RoPE rotation angle, $\frac{n}{\beta^i}$, is employed. When tokens are located at positions where $n \geq \hat{n}$, the scaling factor is then incorporated.

Assuming a desired context window length of $L'$, our task is to determine the set of optimal rescale factors, denoted as ($\mathbb{I}(\hat{\lambda}_{0},\hat{n})$, $\mathbb{I}(\hat{\lambda}_{1},\hat{n})$, ... ,$\mathbb{I}(\hat{\lambda}_{i},\hat{n})$, ...), across the $1^{st}$ to $d^{th}$ dimensions of the RoPE mechanism. The aim is for the modified target $\text{LLM}$—incorporating these adjusted $\text{RoPE}$ factors—to minimize next token prediction loss, $\mathcal{L}$ (measured by perplexity), when processing input sequences $\mathbf{X}$ of token length $L'$.

\subsection{Searching the Non-uniform Position Interpolation}

To address the problem formulated in Eq.~\ref{eq:problem}, we present a straightforward yet robust approach that identifies optimal RoPE rescale factors, thereby harnessing the multidimensional and non-uniform aspects present in position embeddings.

\textbf{Search space}. We construct an expansive search space that accounts for both forms of non-uniformity. Table~\ref{tbl:searchspace} details the structure of this search space. In particular, the method assigns a unique rescale factor to each dimension within the RoPE embedding, allowing for individualized optimization. For the sake of simplicity in designing the search process, we perform the search over $\lambda_i$ and $\hat{n}$ rather than directly searching for $\mathbb{I}(\hat{\lambda}_{i},\hat{n})$, with $\hat{\lambda}_{i}=1/\lambda_i$. As indicated in Table~\ref{tbl:searchspace}, the allowable range for $\lambda_i$ starts at 1.0 (corresponding to direct extrapolation) and extends up to $s\times1.25$ (representing a broader interpolation than Position Interpolation (PI)), incremented in steps of 0.01. Here, $s$ denotes the intended extension ratio of the context window.

The parameter $\hat{n}$ determines how many of the initial token positions preserve the original RoPE embedding, foregoing position interpolation. In practice, we explore values of $\hat{n}$ from the set \{0, 1, 2, 4, 8, 12, 16, 20, 24, 28, 32, 64, 128, 256\}. Setting $\hat{n}=0$ implies that every token position utilizes the optimized rescale factors.

\textbf{Evolution-based search}. The search space outlined in Table~\ref{tbl:searchspace} encompasses a vast array of positional interpolation options, making thorough exploration highly challenging. For instance, an extension factor of $s=4\times$ generates $400^{128/2}\times14$ = 4$\times10^{167}$ possible configurations. As the extension ratio increases, the search space grows exponentially. To efficiently navigate this expansive space, we employ evolution search~\cite{guo2020single} and incorporate two optimization strategies designed to significantly enhance the search process. The complete search framework is described in Algorithm~\ref{alg:search}.

\textit{Optimized initial population generation}. In place of a fully random initialization for the population of $P$ rescale factors, we explicitly include the three RoPE rescale factors associated with PI, NTK, and YaRN within the starting population. For the other $P$-3 individuals, we introduce variation by randomly mutating these three rescale factors with a probability $p$.

\textit{Monotonically non-decreasing constraint}. Once the initial population is formed, we evaluate the perplexity of each candidate using the LLM. To do this, the respective RoPE rescale factors are applied to the target LLM, and we then assess the perplexity of the input $\mathbf{X}$. The best-performing $k$ individuals are selected as parents for the subsequent evolutionary steps. Nonetheless, due to the immense size of the search space, basic mutation and crossover operations can result in suboptimal candidates, incurring excessive perplexity evaluations. This challenge becomes more pronounced for large $L'$, as each perplexity assessment requires a significant amount of inference time.

To tackle this issue, we enforce a constraint of non-decreasing monotonicity on the rescaled RoPE factors being sampled, such that $\lambda_i \le \lambda_{i+1}$. Only those RoPE modifications that meet this requirement are considered when evaluating perplexity in LLMs, which helps to drastically narrow the search space. In detail, our approach stipulates that $\lambda_i$ must monotonically increase across the RoPE dimensions (i.e., for $i=0,\ldots,63$). This design choice is informed by NTK theory~\cite{jacot2018neural,tancik2020fourier,ntk}, which indicates that in lower dimensions—corresponding to higher frequencies—less interpolation is needed (implying a smaller $\lambda_i$), while in higher dimensions—associated with lower frequencies—greater interpolation is permissible (requiring a larger $\lambda_i$).

\begin{algorithm}[t]
	\small
	\caption{The search algorithm}
	\textbf{Input:} target LLM, input samples $\mathbf{X}$, population size $P$, mutation size $N_1$, crossover size $N_2$, maximum iterations $\mathcal{T}$, mutation probability $p$ \\

	\begin{algorithmic}[1]
		\label{alg:search}
		\STATE Top-k $\gets \phi$;	
		\STATE $\text{P}_0 \gets$ \textit{Initialize\_population\_with\_optimization}($P$, $\mathbf{X}$, $p$);
		\FOR{$i = 1$ to $\mathcal{T}$}
		\STATE \textit{Compute\_perplexity}(LLM, $\text{P}_{i-1}$, $\mathbf{X}$);
		\STATE Top-k $\gets$ \textit{Update\_Topk}(Top-k, $\text{P}_{i-1}$);
		\STATE $\text{P}_{mutation} \gets$ \textit{Mutation\_with\_mono\_constraint}(Top-k, $N_1$, $p$);
		\STATE $\text{P}_{crossover} \gets$ \textit{Crossover\_with\_mono\_constraint}(Top-k, $N_2$);
		\STATE $\text{P}_i = \text{P}_{mutation} \cup \text{P}_{crossover} \cup \text{Top-k}$;
		\ENDFOR
		\STATE Return the candidate with minimal perplexity from Top-k;
	\end{algorithmic}
\end{algorithm}

\textbf{8$\times$ extension without fine-tuning}. Utilizing evolutionary search, our approach efficiently determines optimal non-uniform RoPE rescale factors, safeguarding crucial dimensions and positional information to reduce the loss from interpolation. As illustrated in Fig.\ref{fig:non-ft}, this method enables LLaMA2’s context window to expand from 4k to 32k tokens without requiring any fine-tuning. By comparison, prior techniques including PI, as well as non-uniform NTK and YaRN, lead to significant increases in perplexity after only a 2$\times$ extension.

\subsection{Extending LLM Context Window to 2048K}
\label{sec:extendto2048k}

\textbf{Progressive extension to 2048k}. In this section, we present our approach for expanding the context window of pre-trained LLMs far beyond the standard 4k, reaching up to 2048k. Our results show that non-uniform positional interpolation enables us to achieve an 8$\times$ increase in context length without any need for fine-tuning. However, when aiming for more substantial extensions, such as up to 512$\times$, fine-tuning becomes indispensable. One potential strategy involves optimizing RoPE rescaling parameters at the target 2048k length, followed by fine-tuning. Yet, this approach encounters practical limitations due to the immense computational cost involved. Additionally, our observations indicate that effectively fine-tuning LLMs at such a high extension ratio proves to be quite difficult (refer to Appendix).

Fortunately, LongRoPE demonstrates efficacy with both baseline and fine-tuned extended LLMs. Thus, we propose an efficient, incremental approach that attains the desired 2048k context window using only 1k fine-tuning steps, all performed within a training length of 256k.

$\Diamond$ \textit{Exploring LongRoPE for scaling pre-trained LLMs to 256k context windows}. Using LLaMA2 as a case study, we experiment with expanding the context window to sizes of 128k and 256k. This process involves increasing the extension factor to 32$\times$ and 64$\times$ at the corresponding stages.

$\Diamond$ \textit{Fine-tuning to 256k}. Following pre-training, we adapt the LLM to handle a 256k context window through fine-tuning. Our approach begins by fine-tuning LLaMA2 with RoPE rescaled factors configured for 128k over the course of 400 steps. Afterward, we update the model by switching the RoPE rescaled factors to those for 256k using the resulting checkpoint and proceed with 600 additional fine-tuning steps. This staged procedure is demonstrated to be more efficient than immediately fine-tuning for 256k.

$\Diamond$ \textit{Expanding fine-tuned extended LLM to 2048k via LongRoPE search}. In the final step, an additional search is conducted using the LLM that has already been fine-tuned on sequences of 256k length. As a result, the model achieves an exceptionally large context window of 2048k, all without necessitating further fine-tuning. The ultimate expansion factor reaches $512\times$.

\textbf{Performance degradation in original context window}.  When expanding the context window to a substantial 2048k length, we observe that the model's effectiveness within its initial context range decreases. This phenomenon, documented in positional interpolation research~\cite{pi}, occurs because the positional embeddings corresponding to the original window are compressed into a relatively small segment of the embedding space, impairing the language model's accuracy. In scenarios with a 512$\times$ expansion ratio, the original 4k context positions are especially densely clustered.

To address this issue, we conduct an additional evolution search on the expanded LLM, specifically tuning the RoPE rescale factors for smaller context sizes (such as 4k and 8k). Since shorter contexts need less positional interpolation, we constrain the upper limit for the searched $\lambda$. At inference time, the LLM adaptively modifies the relevant RoPE rescale factors.

 \section{Experiments}

\label{sec:exp}
\subsection{Setup}
\textbf{Evaluation Tasks and models}. LongRoPE is integrated with both LLaMA2-7B and Mistral-7B, and their effectiveness is assessed using three criteria: (1) perplexity scores on lengthy documents to gauge how extended-context LLMs manage long-range dependencies; (2) performance on the Passkey retrieval task, which evaluates how accurately the model can extract a target passkey hidden among abundant distractor text; and (3) results on established LLM benchmark datasets, restricted to a 4096 token context window.

\textbf{Fine-tuning}. For LLaMA2, we adopt a learning rate of 2e-5 with a linear decay schedule, utilizing a global batch size of 32. The model is initially fine-tuned over 400 steps on the Redpajama~\cite{redpajama} dataset, divided into 128k-token segments and delimited with BOS and EOS tokens. Building upon the resulting checkpoint, an additional 600 steps of training extend the context window to 256k tokens. Training with the 128k context window employs 8 A100 GPUs through the distributed training system~\cite{cube}, while the 256k variant requires 16 A100 GPUs. For Mistral, a fixed learning rate of 1e-6 is used alongside a global batch size of 64. The procedure for both 128k and 256k models follows the YaRN~\cite{yarn} setup: 400 training steps on Together Computer's Long-Data Collections~\cite{mistraldata} with a 16k sequence length, leveraging 4 A100 GPUs for training.

\textbf{Search}. For target window sizes up to 256k, we employ the following hyperparameters: $P$=64, $N_1$=$N_2$=16, $p$=0.3, and $\mathcal{T}$=40, with the top-32 candidates from each generation chosen for mutation or crossover. Perplexity is determined using five randomly chosen samples from the PG19 validation set, each meeting the minimum target context length. For window sizes exceeding 512k, the population size as well as the numbers used for mutation and crossover are reduced by half. In this case, perplexity is evaluated on three randomly sampled entries from the Pile-Books3~\cite{pile} validation set.

\textbf{Baselines}. For the 2048k context window experiments, we utilized fine-tuned models with context lengths of 128k and 256k, resulting in LongRoPE-2048k (ft=128k) for LLaMA2 and LongRoPE-2048k (ft=256k) for Mistral. These four models are evaluated against leading context window extension baselines. In particular, we consider open-source LLMs that have undergone fine-tuning following positional interpolation techniques such as PI, NTK, and YaRN. The baselines include Together-32k~\cite{together}, Code LLaMA~\cite{codellama}, LongLoRA-full-FT-100k~\cite{longlora}, as well as YaRN-LLaMA and YaRN-Mistral~\cite{yarn}.

\subsection{Main Results}

\label{sec:mainresult}
\textbf{Long sequence language modeling up to 256k}. To start, we assess performance in language modeling for sequences as long as 256k, comparing our method with leading extended LLMs. Two datasets, Proof-pile~\cite{pg19} and the test split of PG19~\cite{pile}, are utilized to highlight the broad applicability of our approach. Perplexity is measured across a range of context window lengths, employing a sliding window of size 256. For PG19, evaluation is performed on the complete test split, comprising 100 documents. For Proof-pile, in alignment with YaRN~\cite{yarn}, we randomly select 10 samples, each with a minimum sequence length of 128k.

Table~\ref{tbl:proof-pile} and Table~\ref{tbl:pg19} present a comparison of the perplexity results for LLaMA2 and Mistral models augmented with various interpolation techniques, evaluated on Proof-pile and PG19, respectively. From these results, two primary insights emerge: \textbf{(1)} the extended models consistently achieve lower perplexity as the evaluation length increases from 4k up to 256k, demonstrating effective utilization of longer contextual information. \textbf{(2)} Notably, even when operating with a context window expanded by a factor of 16—a situation that usually hampers performance at shorter ranges—our LongRoPE-2048k variants surpass leading baseline models within the 256k context length range.

\textbf{Modeling language on sequences exceeding 2000k tokens}. To assess performance on exceptionally lengthy documents, we utilize the Books3~\cite{pile} dataset. For efficient evaluation, we randomly choose 20 books with lengths greater than 2048k, applying a sliding window of 256k tokens.

As indicated in Table~\ref{tbl:books3}, LongRoPE is able to increase the context window of both LLaMA2-7B and Mistral-7B models up to 2048k, all the while maintaining perplexity levels that are on par with, or better than, baseline approaches across the shorter range of 8k to 128k. Substantial differences in performance between the 2048k variants of LLaMA2 and Mistral are apparent. While Mistral exhibits superior results compared to baselines at shorter context lengths, its perplexity rises above 7 when the context extends beyond 256k. LLaMA2 demonstrates behavior consistent with anticipated trends, as perplexity continues to improve with extended contexts, only displaying minor increases at 1024k and 2048k. Additionally, on LLaMA2, LongRoPE-2048k achieves enhanced results when fine-tuned at 256k rather than 128k, a consequence of employing a reduced secondary extension ratio (i.e., 8$\times$ versus 16$\times$). In contrast, for Mistral, the model yields improved outcomes with a fine-tuning window of 128k. This primarily stems from our adoption of the YaRN configuration for Mistral's 128k and 256k fine-tuning, which utilizes a 16k training length and subsequently influences Mistral's capacity for post-fine-tuning context window extension.

\textbf{Passkey retrieval}. Next, we examine how the usable context window size influences performance on generation tasks. Adopting the synthetic passkey retrieval evaluation outlined by ~\cite{passkey}, the model must extract a random five-digit passkey that is concealed within a lengthy document. The structure of the prompt can be found in the appendix. For this experiment, we conduct 10 rounds of the retrieval task, each time embedding the passkey at a randomly selected position uniformly spread throughout the evaluation context window.

Fig.~\ref{fig:passkey} presents a comparison of retrieval accuracy against baseline approaches. The performance of existing models quickly deteriorates, with accuracy falling to zero when the sequence length exceeds 128k tokens. In stark contrast, LongRoPE-LLaMA2-2048k (ft=256k) consistently achieves high retrieval accuracy ($\ge$90\%) across a broad range of context lengths, from 4k up to 2048k tokens, even under the demanding condition of recovering a passkey from sequences containing millions of tokens. Meanwhile, LongRoPE-Mistral-2048k (ft=128k) sustains perfect accuracy (100\%) until 1800k tokens and then decreases to 60\% at 2048k, which is consistent with the pattern observed in Table~\ref{tbl:books3} where perplexity experiences a slight increase at the 2048k context length.

\textbf{Standard benchmarks in the native context window}. LongRoPE-2048k models are assessed within their original context window through the Hugging Face Open LLM Leaderboard~\cite{huggingfaceboard}, utilizing both zero-shot and few-shot paradigms. The evaluation includes the 25-shot ARC-Challenge~\cite{allenai:arc}, 10-shot HellaSwag~\cite{zellers2019hellaswag}, 5-shot MMLU~\cite{mmlu}, and 0-shot TruthfulQA~\cite{lin2021truthfulqa} tasks.

As indicated in Table~\ref{tbl:commonsense}, our models attain results similar to those observed on the initial benchmark, which was developed with a shorter context window in mind, and surpass the original Mistral on TruthfulQA by a margin of +0.5\%. LongRoPE-LLaMA2-2048k, after being fine-tuned on 256k tokens, exhibits a marginally higher drop in performance, yet continues to deliver acceptable outcomes on the majority of tasks.

\subsection{Ablation Results}

\textbf{Effectiveness of the second positional interpolation}. Our progressive extension approach incorporates a second non-uniform positional interpolation step, applied through our search algorithm to further extend fine-tuned LLMs. To assess this method, we conducted tests on the fine-tuned LLaMA2-256k model, incrementally enlarging its context window to 512k, 1024k, and 2048k tokens via both PI and YaRN techniques. According to the results in Table~\ref{tbl:secondsearch}, our non-uniform positional interpolation method maintains perplexity at a stable level across expansions. In comparison, perplexity rises rapidly as the extension ratio increases for both PI and YaRN approaches.

\textbf{Effectiveness of recovery at shorter context lengths.} To address reductions in performance at shorter context lengths, we re-optimized the RoPE scaling parameters for LongRoPE-2048k using our search procedure. In particular, we imposed tighter constraints on the maximum scale factors during the search process to limit interpolation when evaluating at shorter 4k and 8k contexts. Table~\ref{tbl:shortperformance} presents the perplexity metrics for LongRoPE-LLaMA2-2048k on Proof-pile for 4k and 8k lengths, alongside mean LLM benchmark scores. These findings highlight substantial improvements in performance at short context lengths.

\textbf{Analysis on the two forms of non-uniformities}. To assess the impact of the two non-uniformities on overall performance, we perform an ablation study. Our investigation consists of two primary experiments: (i) for LLaMA2-7B, we extend the context window to 16k and 32k tokens using various strategies—including PI, searching solely across the RoPE dimension, and exploring both forms of non-uniformity; (ii) for our 256k-length fine-tuned LLaMA2 model, we extrapolate to 2048k using analogous approaches. In all cases, we report perplexity scores without any additional fine-tuning. As indicated in Table~\ref{tbl:searchdimension}, adjusting the RoPE dimension to allow for non-uniformity yields a marked decrease in perplexity relative to the linear interpolation used in PI. While introducing non-uniformity over token positions boosts performance for sequence lengths of 16k and 32k, its effect becomes negligible at the extreme 2048k context length, possibly as a consequence of the vast input size. Notably, merely retaining the earliest tokens and omitting interpolation does not offer discernible benefits, suggesting that further investigation is warranted for this approach.

\section{Related Works}

Beyond techniques that rely on position interpolation, this section examines alternative methodologies.

\textbf{Retrieval-based} approaches leverage external memory components to store distant contextual information and employ retrieval mechanisms to fetch pertinent documents during inference~\cite{longllama,wang2023augmenting,borgeaud2022improving}. Such systems often require substantial architectural alterations to large language models. By contrast, our method remains comparatively simple, introducing only modest changes to positional embeddings. Furthermore, our approach accommodates a broader range of long-context applications extending past retrieval tasks, including lengthy document summarization and few-shot learning.

\textbf{Attention-based context window extensions}. In addition to positional embedding interpolation, certain studies have managed to extend the input context window—while maintaining the LLM’s native context length—through modifications of the attention mechanism~\cite{han2023lminfinite, streamingllm,parallelwindow}. Central to this approach is addressing the attention explosion problem induced by novel positional encodings via the introduction of innovative attention masks. Such approaches are complementary to positional interpolation techniques.

\textbf{Fine-tuning based approaches} are concerned with optimizing the adaptation of pre-trained LLMs that utilize altered position embeddings to accommodate longer input sequences. Methods such as Code LLaMA~\cite{codellama}, LLaMA2 Long~\cite{xiong2023effective}, and ScaledRoPE~\cite{liu2023scaling} adopt a substantially larger base for RoPE and proceed to fine-tune the models for the desired sequence length. The technique we introduce enables support for a wide range of target lengths, surpassing 2M tokens. In recent developments, due to the considerable computational cost of fine-tuning models on extremely long contexts (greater than 128k tokens), solutions like LongLoRA~\cite{longlora} and PoSE~\cite{zhu2023pose} have emerged to alleviate such resource requirements. Our approach is complementary to these efficient fine-tuning methods.

\section{Conclusion}

In this study, we introduce LongRoPE, a technique that significantly increases the context length of LLMs to a groundbreaking 2048k, all while preserving their performance within the original, shorter context windows. Our approach leverages two types of non-uniformities in the RoPE positional encoding, utilizing an efficient evolutionary search process. This method yields dual advantages: it facilitates optimal initialization for subsequent fine-tuning and allows an 8$\times$ expansion of the context window even in the absence of fine-tuning. Furthermore, we put forward a progressive extension framework, wherein LLMs fine-tuned with a 256k context length are utilized to ultimately achieve a 2048k window size without the need for additional fine-tuning steps. Comprehensive experimental evaluations confirm the robustness and effectiveness of LongRoPE. We anticipate that the LongRoPE-2048k models will support a variety of new applications requiring extensive context and stimulate further investigation in this area.

\section*{Broader Impacts} The objective of this research is to contribute to progress in the field of Machine Learning. While our work may carry a range of possible implications for society, we do not identify any particular issues that require special mention at this time.

	\balance

\end{document}